% Riemann--Roch and Its Consequences  (source §6; restructured: the dimension
% counts that justify 1, wp, wp' are derived ELEMENTARILY from Chapter 4 plus the
% explicit Weierstrass functions, so the main line is self-contained; Riemann--Roch
% is kept as the general theorem (stated, cited) and is genuinely needed only for
% the converse of Abel / Pic^0(T) =~ T, which is off the critical path.)
\chapter{Riemann--Roch and Its Consequences}
\label{ch:riemann_roch}

We can now answer the question left open in \cref{ch:elliptic_liouville}: how
small can the pole divisor of a nonconstant elliptic function be? The answer is a
dimension count, and on the torus that count is \emph{elementary} --- it follows
from the Liouville theorems of \cref{ch:elliptic_liouville} together with the
explicit functions built in \cref{ch:weierstrass_wp}, with no appeal to a deep
theorem. We derive it that way first, because it is what justifies the functions
\(1, \wp, \wp'\) and carries the main argument. We then state the
\emph{Riemann--Roch} theorem in general, for context. The one auxiliary result it
might have supplied --- the converse of Abel's law, \(\Pic^0(T) \cong T\) --- is
instead proved constructively in \cref{ch:weierstrass_wp} with the
\(\sigma\)-function, so Riemann--Roch is needed nowhere and appears only as the
general theorem behind the elementary counts.

\section{The Spaces \texorpdfstring{$L(D)$}{L(D)}}
\label{sec:rr-spaces}

We order divisors by \(D \ge 0\) iff every coefficient of \(D\) is
nonnegative (an \emph{effective} divisor), and \(D \ge D'\) iff \(D - D' \ge 0\).

\begin{definition}
\label{def:rr-LD}
For a divisor \(D\) on a compact Riemann surface \(X\),
\[
  L(D) = \{ f \in \M(X) : \divisor(f) + D \ge 0 \} \cup \{0\},
  \qquad \ell(D) = \dim_\C L(D) .
\]
The condition \(\divisor(f) + D \ge 0\) says \(f\) may have a pole of order at
most \(n_p\) where \(D = \sum n_p\,p\) has \(n_p > 0\), and must have a zero of
order at least \(-n_p\) where \(n_p < 0\). For \(D = nO\) with \(n \ge 0\),
\(L(nO)\) is the space of functions whose only pole is at \(O\), of order at most
\(n\).
\end{definition}

\begin{proposition}[Elementary properties of \texorpdfstring{$L(D)$}{L(D)}]
\label{prop:rr-LD-props}
Let \(D, D'\) be divisors on a compact Riemann surface \(X\).
\begin{enumerate}
\item \(L(D)\) is a \(\C\)-vector space, and \(\ell(0) = 1\) with \(L(0) = \C\).
\item If \(D \le D'\) then \(L(D) \subseteq L(D')\) and
  \(\ell(D') - \ell(D) \le \deg(D' - D)\).
\item If \(\deg D < 0\) then \(L(D) = 0\), so \(\ell(D) = 0\).
\item Linearly equivalent divisors have isomorphic spaces: \(D \sim D'\) implies
  \(L(D) \cong L(D')\), so \(\ell\) depends only on the class \([D]\).
\end{enumerate}
\end{proposition}

\begin{proof}
(1) If \(\divisor(f) + D \ge 0\) and \(\divisor(g) + D \ge 0\) then for
\(a, b \in \C\), \(\ord_p(af+bg) \ge \min(\ord_p f, \ord_p g) \ge -n_p\)
(\cref{ex:div-ord-welldef}), so \(af + bg \in L(D)\). For \(D = 0\), \(L(0)\) is
the holomorphic functions on \(X\), which are the constants by
\cref{thm:ell-const} (compactness), so \(\ell(0) = 1\).

(2) The inclusion is immediate from \(\divisor(f) + D \le \divisor(f) + D'\).
Writing \(D' = D + \sum_j p_j\) with \(\deg(D'-D)\) points (with multiplicity),
each added point relaxes the pole bound at \(p_j\) by one, and the map sending
\(f\) to its Laurent coefficient of that extra allowed order is linear; its
kernel is the smaller space, so the dimension grows by at most one per point.

(3) A nonzero \(f \in L(D)\) has \(\divisor(f) \ge -D\), so
\(0 = \deg\divisor(f) \ge -\deg D\) by \cref{thm:div-degzero}; thus
\(\deg D \ge 0\). Contrapositively \(\deg D < 0\) forces \(L(D) = 0\).

(4) If \(D' = D + \divisor(h)\), then \(f \mapsto f/h\) is a linear isomorphism
\(L(D) \to L(D')\), since \(\divisor(f/h) + D' = \divisor(f) + D \ge 0\).
\end{proof}

\section{The Dimension Counts, Elementarily}
\label{sec:rr-counts}

The counts that the rest of the book uses --- \(\ell(O) = 1\), \(\ell([a]) = 1\),
and \(\ell(nO) = n\) --- need nothing beyond \cref{ch:elliptic_liouville} and the
explicit functions of \cref{ch:weierstrass_wp}.

\begin{proposition}[Dimension counts on the torus]
\label{prop:rr-counts}
On \(T = \C/\Lambda\):
\begin{enumerate}
\item \(\ell(O) = 1\), and \(\ell([a]) = 1\) for every single point \([a]\);
\item \(\ell(nO) \le n\) for all \(n \ge 1\);
\item \(\ell(nO) = n\) for all \(n \ge 1\).
\end{enumerate}
\end{proposition}

\begin{proof}
(1) A nonzero \(f \in L([a])\) has at most a simple pole at \([a]\) and no other
poles, so it is an elliptic function of order at most one. By
\cref{cor:ell-order2} there is no elliptic function of order exactly one, so \(f\)
is holomorphic, hence constant (\cref{thm:ell-const}); thus \(\ell([a]) = 1\).
The case \([a] = O\) is the same.

(2) By the point lemma \(\ell(D + p) - \ell(D) \in \{0,1\}\)
(\cref{ex:rr-point-lemma}), applied \(n-1\) times starting from \(\ell(O) = 1\),
\[
  \ell(nO) \le \ell(O) + (n-1) = n .
\]

(3) For the matching lower bound we exhibit \(n\) independent elements of
\(L(nO)\). \Cref{ch:weierstrass_wp} constructs the elliptic functions \(\wp\)
(a double pole at \(O\), order \(2\)) and \(\wp'\) (a triple pole, order \(3\)).
The functions
\[
  1, \quad \wp^{a}\ (a \ge 1),\quad \wp^{a}\wp'\ (a \ge 0)
\]
have poles only at \(O\), of orders \(0\) and \(2a\) and \(2a+3\) respectively, so
they realize every pole order in \(\{0, 2, 3, 4, \dots, n\}\) --- that is \(n\)
distinct orders, hence \(n\) linearly independent members of \(L(nO)\). Therefore
\(\ell(nO) \ge n\), and with (2), \(\ell(nO) = n\).
\end{proof}

\begin{remark}
\label{rem:rr-elementary}
The proof used only the Liouville theorems of \cref{ch:elliptic_liouville}
(\cref{cor:ell-order2,thm:ell-const}), the degree-zero theorem
(\cref{thm:div-degzero}, via the point lemma), and the existence of \(\wp, \wp'\)
from \cref{ch:weierstrass_wp}. No Riemann--Roch is needed. The same argument gives
\(\ell([a]) = 1\) for any point, which is all the converse of Abel will require of
these counts.
\end{remark}

\section{The Canonical Divisor of a Torus}
\label{sec:rr-canonical}

For the general theorem we need the \emph{canonical divisor} \(K\), the divisor of
any nonzero meromorphic \(1\)-form. On a compact surface of genus \(g\) all such
forms are linearly equivalent and \(\deg K = 2g - 2\). For the torus the canonical
class is trivial.

\begin{proposition}
\label{prop:rr-canonical}
On \(T = \C/\Lambda\) the holomorphic \(1\)-form \(dz\) is translation-invariant
and nowhere vanishing, so
\[
  K = \divisor(dz) = 0 .
\]
Hence the canonical class is trivial, \(\deg K = 0\), and the space of
holomorphic \(1\)-forms is one-dimensional, \(\Omega(T) = \C\,dz\); in particular
\(g = \dim\Omega(T) = 1\).
\end{proposition}

\begin{proof}
Since \(d(z+\lambda) = dz\) for every \(\lambda \in \Lambda\), the form \(dz\)
descends to \(T\); in the charts of \cref{thm:lat-riemann-surface} it is simply
\(dz\), which has no zeros or poles, so \(\divisor(dz) = 0\). Any holomorphic
\(1\)-form is \(f\,dz\) with \(f\) holomorphic on \(T\), hence constant
(\cref{thm:ell-const}); thus \(\Omega(T) = \C\,dz\) is one-dimensional. The
relation \(\deg K = 2g - 2\) with \(\deg K = 0\) gives \(g = 1\), agreeing with
\cref{prop:tg-euler}.
\end{proof}

\section{Riemann--Roch: The General Theorem}
\label{sec:rr-theorem}

For surfaces of any genus the dimension count is governed by one theorem, which we
state for context. On the torus it reproduces the counts of \cref{prop:rr-counts}
and extends them to \emph{arbitrary} divisors; we include it for context only.
Even the converse of the next section --- the one result beyond the elementary
counts --- is obtained constructively in \cref{ch:weierstrass_wp} via the
\(\sigma\)-function, so nothing in the book depends on this theorem.

\begin{theorem}[Riemann--Roch]
\label{thm:rr}
Let \(X\) be a compact Riemann surface of genus \(g\) with canonical divisor
\(K\). For every divisor \(D\),
\[
  \ell(D) - \ell(K - D) = \deg D + 1 - g .
\]
\end{theorem}

A proof belongs to a course on Riemann surfaces; see, e.g., Miranda
\cite{miranda} or Forster \cite{forster}. The term \(\ell(K-D)\) is the
\emph{index of speciality}; by
Serre duality it equals the dimension of holomorphic \(1\)-forms whose divisor is
\(\ge D\). On the torus \(g = 1\) and \(K \sim 0\) (\cref{prop:rr-canonical}), so
\(\ell(K - D) = \ell(-D)\) (\cref{prop:rr-LD-props}(4)) and Riemann--Roch reads
\[
  \boxed{\ \ell(D) - \ell(-D) = \deg D\ } .
\]

\begin{proposition}[Dimensions for arbitrary \(D\), from Riemann--Roch]
\label{prop:rr-dims}
On \(T\), for any divisor \(D\): if \(\deg D < 0\) then \(\ell(D) = 0\); if
\(\deg D > 0\) then \(\ell(D) = \deg D\); and if \(\deg D = 0\) then
\(\ell(D) = 1\) for \(D \sim 0\) and \(0\) otherwise. In particular
\(\ell(D) = \deg D\) for every divisor of degree \(\ge 1\), whether or not \(D\)
is supported at \(O\).
\end{proposition}

\begin{proof}
For \(\deg D < 0\), \cref{prop:rr-LD-props}(3). For \(\deg D > 0\), \(-D\) has
negative degree so \(\ell(-D) = 0\) and the boxed identity gives
\(\ell(D) = \deg D\). For \(\deg D = 0\), the identity gives
\(\ell(D) = \ell(-D)\); a nonzero \(f \in L(D)\) has \(\divisor(f) \ge -D\) of
degree \(0\), forcing \(\divisor(f) = -D\), i.e.\ \(D \sim 0\), so
\(\ell(D) \le 1\) with equality iff \(D \sim 0\).
\end{proof}

For divisors supported at \(O\) this only restates \cref{prop:rr-counts}; its new
content, the value for \emph{arbitrary} degree-one divisors, is what
Riemann--Roch would contribute --- but the next section reaches the same end
constructively, so we do not rely on it.

\section{The Converse of Abel: Completing \texorpdfstring{$\Pic^0(T) \cong T$}{Pic0(T) = T}}
\label{sec:rr-pic-completion}

We can now finish the injectivity left open in \cref{thm:div-pic-iso}.

\begin{corollary}
\label{cor:rr-degree-one}
Every divisor of degree \(1\) on \(T\) is linearly equivalent to a unique point
\([a]\). Consequently every degree-zero divisor is linearly equivalent to a unique
\([a] - [0]\), and the Abel--Jacobi map \(\bar\smap : \Pic^0(T) \to T\) of
\cref{thm:div-pic-iso} is an isomorphism.
\end{corollary}

\begin{proof}
The key input is the constructive converse of Abel, \cref{thm:wp-converse}: a
degree-zero divisor with Abel sum zero is principal. \emph{Injectivity:} if
\(\smap(D) = 0\) for \(\deg D = 0\), then \(D\) is principal, so its class in
\(\Pic^0(T)\) is trivial; with surjectivity from \cref{thm:div-pic-iso},
\(\bar\smap\) is an isomorphism. \emph{Degree-one statement:} for \(\deg D = 1\)
set \(a := \smap(D - [0]) \in T\); then \(D - [a]\) has degree \(0\) and Abel sum
\(\smap(D-[0]) - a = 0\), so by \cref{thm:wp-converse} \(D - [a]\) is principal and
\(D \sim [a]\). Uniqueness: if \([a] \sim [a']\) then \([a] - [a']\) is principal of
degree \(0\), so by \cref{ex:div-distinct-points} (the forward direction of Abel)
\([a] = [a']\).
\end{proof}

\begin{remark}[Where existence is genuinely used]
\label{rem:rr-existence}
This corollary is the one result needing more than the elementary counts: the
converse of Abel, that a balanced degree-zero divisor is principal. We do
\emph{not} import it from Riemann--Roch --- \cref{thm:wp-converse} supplies it
constructively, building the required elliptic function explicitly from the
Weierstrass \(\sigma\)-function. So even this is proved from within the book.
Moreover it is \emph{not} on the critical path to the main theorem:
\cref{thm:tr-main} uses only the necessity direction of Abel (\cref{thm:ell-abel})
and the elementary counts of \cref{prop:rr-counts}. The isomorphism
\(\Pic^0(T) \cong T\) is the algebraic \emph{view} of the group law, proved in full
but not a step in the chord--tangent argument.
\end{remark}

\section{The Functions Forced}
\label{sec:rr-functions}

By \cref{prop:rr-counts}, \(\ell(nO) = n\) for \(n \ge 1\): the dimensions climb by
exactly one at each step, so each \(L((n{+}1)O)\) contains a function with a pole
at \(O\) of order exactly \(n+1\) not already present.

\begin{itemize}
\item \(\ell(O) = 1\): \(L(O) = \C\). There is \emph{no} nonconstant function with
  a single simple pole --- the count restates \cref{cor:ell-order2}.
\item \(\ell(2O) = 2\): beyond the constants there is one new function, with a
  double pole at \(O\). After normalization this is the Weierstrass function
  \(\wp\).
\item \(\ell(3O) = 3\): one further new function, with a triple pole at \(O\),
  namely \(\wp'\). A basis of \(L(3O)\) is \(\{1, \wp, \wp'\}\).
\end{itemize}

\begin{center}
\begin{tabular}{c|c|c|l}
Divisor & \(\ell\) & Typical basis & Meaning\\
\hline
\(O\)  & \(1\) & \(1\)            & no simple-pole function\\
\(2O\) & \(2\) & \(1, \wp\)       & first nonconstant function\\
\(3O\) & \(3\) & \(1, \wp, \wp'\) & projective embedding data
\end{tabular}
\end{center}

\begin{keyidea}{Why \(\wp\) and \(\wp'\) are inevitable}
The dimension counts (\cref{prop:rr-counts}) force the first nonconstant function
on the torus to appear with a double pole, and its derivative then has a triple
pole. The first three natural functions are therefore
\[
  1, \qquad \wp, \qquad \wp' ,
\]
and these are exactly the coordinates of the map that embeds the torus as a plane
cubic in \cref{ch:embedding}.
\end{keyidea}

\begin{remark}[The cubic is forced too]
\label{rem:rr-cubic-preview}
Push the count one step further. In \(L(6O)\), which has dimension \(\ell(6O) =
6\), the seven functions
\[
  1,\ \wp,\ \wp^2,\ \wp^3,\ \wp',\ \wp\wp',\ \wp'^2
\]
all lie (their pole orders at \(O\) are \(0,2,4,6,3,5,6\), each \(\le 6\)). Seven
vectors in a six-dimensional space are linearly dependent, so they satisfy a
linear relation. Matching the order-six poles forces \(\wp'^2\) and \(\wp^3\) to
appear with nonzero coefficients: the relation is the Weierstrass cubic
\(\wp'^2 = 4\wp^3 - g_2\wp - g_3\), derived explicitly in
\cref{ch:weierstrass_wp}.
\end{remark}

\section{Summary}
\label{sec:rr-summary}

The space \(L(D)\) measures functions with poles bounded by \(D\)
(\cref{prop:rr-LD-props}). On the torus the counts that matter --- \(\ell(O) = 1\),
\(\ell([a]) = 1\), and \(\ell(nO) = n\) --- are \emph{elementary}, following from
the Liouville theorems and the explicit \(\wp, \wp'\) with no Riemann--Roch
(\cref{prop:rr-counts}); they force the functions \(1, \wp, \wp'\)
(\cref{sec:rr-functions}) and already predict the cubic
(\cref{rem:rr-cubic-preview}). Riemann--Roch (\cref{thm:rr}) is the general theorem
behind them; on the torus it collapses to \(\ell(D) - \ell(-D) = \deg D\)
(\cref{prop:rr-dims}). Even the converse of Abel, giving \(\Pic^0(T) \cong T\)
(\cref{cor:rr-degree-one}), is proved constructively in \cref{ch:weierstrass_wp}
via the \(\sigma\)-function (\cref{rem:rr-existence}), so Riemann--Roch is needed
nowhere and serves only as the general theorem. The next chapter constructs
\(\wp\) --- and \(\sigma\) --- and verifies all of this by hand.

\section*{Exercises}
\addcontentsline{toc}{section}{Exercises}

\begin{exercise}
\label{ex:rr-point-lemma}
Prove the ``point lemma'': for any divisor \(D\) and point \(p\),
\(\ell(D + p) - \ell(D) \in \{0, 1\}\). (Map \(f \in L(D+p)\) to the Laurent
coefficient of order \(-n_p - 1\) at \(p\); identify the kernel.)
\end{exercise}

\begin{exercise}
\label{ex:rr-neg-degree}
Without invoking Riemann--Roch, show directly from \cref{thm:div-degzero} that
\(\ell(D) = 0\) whenever \(\deg D < 0\), and that \(\ell(D) \le \deg D + 1\) for
all \(D\). (The second is the ``Riemann inequality'' half of the theorem.)
\end{exercise}

\begin{exercise}
\label{ex:rr-degree-one-unique}
Show that on \(T\) a divisor of degree one is effective if and only if it is a
single point, and combine this with \cref{prop:rr-counts} to give an independent
proof that the only effective divisor in the class of \([a]\) is \([a]\) itself.
\end{exercise}

\begin{exercise}
\label{ex:rr-canonical-deg}
Using \cref{prop:rr-canonical} and Riemann--Roch with \(D = K\), verify the two
consistency checks \(\ell(K) = g\) and \(\deg K = 2g - 2\) for the torus
(\(g = 1\)).
\end{exercise}

\begin{exercise}
\label{ex:rr-cubic-count}
Carry out the counting of \cref{rem:rr-cubic-preview} in detail: list the pole
orders at \(O\) of the seven functions
\[
  1,\quad \wp,\quad \wp^2,\quad \wp^3,\quad \wp',\quad \wp\wp',\quad \wp'^2,
\]
explain why all seven lie in \(L(6O)\), and argue that any linear relation among
them must have nonzero coefficients on both \(\wp'^2\) and \(\wp^3\). What goes
wrong if either coefficient is zero?
\end{exercise}
